\section{Model and Problem Statement}\label{sec:model}

\subsection{Private Query-Difference Setting}
\label{subsec:private-query-difference-setting}

We consider a trusted server that holds a private database $D$ and receives, or derives, two related queries $Q_1$ and $Q_2$ over the same database.
The goal is to explain the semantic difference between the corresponding query outputs $Q_1(D)$ and $Q_2(D)$ while preserving the privacy of individuals in $D$.

In the query-refinement setting, $Q_1=Q$ is the original user query and $Q_2=Q'$ is a refined query produced by a refinement mechanism so that $Q'(D)$ satisfies a set of output constraints.
In the consistent-query answering setting, the two outputs correspond to evaluating a query directly on the original database and evaluating it under repair semantics.
In both settings, our framework treats the mechanism that produces the related output as a black box and focuses only on explaining the difference between the two resulting outputs.

\alon{TODO: Add a short paragraph or figure that explicitly describes the trusted-server workflow: the user sends a query to the server, the server holds the private database, derives the related output, and releases differentially private explanations. Replace or update the current placeholder figure accordingly.}

\subsection{Explanation Candidates}
\label{subsec:explanation-candidates}

An explanation in our framework compares the distribution of one attribute across the two query outputs, restricted to the same semantic region of the database.
The semantic region is represented by a predicate view $p(D)$, and the compared distribution is represented by the one-way marginal of an attribute $A$.

\begin{definition}[Histogram Explanation Candidate]
Given a database $D$, two related queries $Q_1,Q_2$, a predicate view $p$, and an attribute $A \in Attr(D)$, a \emph{histogram explanation candidate} is a pair
\[
    e=(p,A).
\]
The candidate induces the two histograms
\[
    H_{Q_1(p(D)),A}
    \qquad\text{and}\qquad
    H_{Q_2(p(D)),A}.
\]
\end{definition}

The predicate view $p(D)$ localizes the explanation to a semantic region of the data, while the attribute $A$ determines the distribution being compared.
A candidate is informative when the two induced histograms differ substantially, since such a difference indicates that the change from $Q_1$ to $Q_2$ affected the selected population with respect to attribute $A$.

\alon{TODO: Add a small running example of a candidate $(p,A)$, e.g., $p=(\texttt{education=Masters})$ and $A=\texttt{occupation}$, and show the two histograms induced by $Q_1$ and $Q_2$.}

\subsection{Ideal Explanation Quality}
\label{sec:explanation-quality}

We evaluate explanation candidates according to four desiderata: \emph{interestingness}, \emph{representativeness}, \emph{coverage}, and \emph{diversity}.
Similar families of criteria are commonly used in explanation and data-exploration systems: interestingness captures whether an explanation highlights a meaningful deviation, while diversity encourages a set of explanations to cover complementary aspects of the data rather than repeatedly showing the same phenomenon \cite{FEDEX}.
Histogram-based explanations and distributional summaries are also used in private explanation frameworks such as DPXPlain \cite{TaoGMR22}.
\alon{TODO: Add a citation to Differentially Private Explanations for Clusters. Also verify whether additional citations are needed for representativeness and coverage.}

For an attribute $A$, let $\operatorname{Dom}(D[A])$ denote the finite set of values, or buckets, used in the histogram of $A$.
For a candidate $e=(p,A)$, the \emph{interestingness} component measures how much the distribution of $A$ changes between $Q_1(p(D))$ and $Q_2(p(D))$:
\[
I(p,A)
=
\sum_{v \in \operatorname{Dom}(D[A])}
\left|
\Pr(t[A]=v \mid t \in Q_1(p(D)))
-
\Pr(t[A]=v \mid t \in Q_2(p(D)))
\right|.
\]
Intuitively, an explanation is interesting if it highlights an attribute whose distribution is substantially affected by the query change.
\alon{TODO: Specify the convention for empty query outputs $Q_i(p(D))$, e.g., whether such candidates are discarded or assigned score $0$.}

The \emph{representativeness} component favors predicate views that are not overly exceptional or idiosyncratic relative to the underlying database.
We measure how close the distribution of $A$ inside the view $p(D)$ is to its distribution in the full database:
\[
R(p,A)
=
1-
\sum_{v \in \operatorname{Dom}(D[A])}
\left|
\Pr(t[A]=v \mid t \in D)
-
\Pr(t[A]=v \mid t \in p(D))
\right|.
\]
This prevents explanations from focusing on small or unusual regions that may exhibit large changes but are not representative of the data.
\alon{TODO: Check normalization. Since TVD includes a factor of $1/2$, decide whether representativeness should be $1-TVD$ or $1-\sum_v |\cdot|$. The current expression follows the previous draft but may range below $0$.}

The \emph{coverage} component favors views that contain sufficient evidence from the compared query outputs.
High coverage makes the explanation statistically more reliable and less likely to be driven by a few tuples:
\[
C(p)
=
|p(D)| \cdot |Q_1(p(D))| \cdot |Q_2(p(D))|.
\]
\alon{TODO: Decide whether the ideal coverage should be the product above, a normalized product, or a sum. The algorithm currently uses a normalized coverage proxy.}

For a single candidate, the ideal quality is
\[
U^*(p,A)
=
I(p,A) \cdot R(p,A) \cdot C(p).
\]
When releasing several explanations, the selected set should also be diverse.
Diversity encourages the explanations to expose different semantic aspects of the change from $Q_1(D)$ to $Q_2(D)$, for example by focusing on different attributes or different views.
\alon{TODO: Add the exact ideal set-level utility, including the diversity term. Possible options are a multiplicative diversity factor, a submodular objective, or a constraint that prevents selecting the same attribute multiple times.}

The ideal objective above captures the explanations we would like to select in the absence of privacy and computational constraints.
However, directly optimizing it is problematic for differentially private selection.
In particular, the score contains conditional distributions and coverage terms that may change significantly when a single tuple is added or removed, especially for small predicate views or small query outputs.
This motivates the low-sensitivity proxy measures introduced in Section~\ref{sec:solution}.

\begin{proposition}[Sensitivity of the ideal quality function]
\label{prop:ideal-utility-sensitivity}
The ideal quality function $U^*$ may have high global sensitivity.
\end{proposition}

\begin{proof}
\alon{TODO: Add proof.}
\end{proof}

\subsection{Problem Definition}
\label{subsec:problem-definition}

We can now state the explanation task using the candidate and quality notions above.
The non-private version asks for high-quality explanations of the difference between two related query outputs.
The private version adds the requirement that the released explanations and histograms satisfy differential privacy with respect to the private database $D$.

\begin{definition}[Query-Difference Explanation Problem]
Given a database $D$, two related queries $Q_1,Q_2$, a set of predicate views $\mathcal{P}$, an attribute set $Attr(D)$, and an integer $k$, return a set $E$ of $k$ histogram explanation candidates $(p,A)$ that maximizes the ideal explanation quality for the semantic difference between $Q_1(D)$ and $Q_2(D)$.
\end{definition}

\begin{definition}[Differentially Private Query-Difference Explanation Problem]
Given a private database $D$, two related queries $Q_1,Q_2$, a set of predicate views $\mathcal{P}$, an attribute set $Attr(D)$, a privacy budget $\epsilon>0$, and an integer $k$, output $k$ histogram-based explanations for the difference between $Q_1(D)$ and $Q_2(D)$ while satisfying $\epsilon$-differential privacy with respect to $D$.
\end{definition}

\alon{TODO: Decide whether the released object should be defined here as noisy difference histograms $\widetilde{\Delta H}_{p,A}$, pairs of noisy histograms $(\hat{H}_{Q_1(p(D)),A},\hat{H}_{Q_2(p(D)),A})$, or both. The current algorithm releases noisy difference histograms.}

\subsection{Challenges}
\label{subsec:model-challenges}

The ideal problem above is difficult for two reasons.
First, the space of possible predicate views is combinatorial in the number of attributes and values, so enumerating all possible semantic regions is infeasible.
Second, the ideal utility function may have high sensitivity, which would require adding large amounts of noise or would make private selection ineffective.
Therefore, the algorithm in the next section restricts the candidate space to predicate-based views and replaces the ideal utility components with low-sensitivity proxy measures.
